% 数学ノート集・共通テンプレート v2（積分論2026年度第1部の版面を参考）
% A5判。章内で定義・定理・例・注意等の番号を共通化。
% コンパイル： platex -interaction=nonstopmode math_note_template_v2.tex (2回)
%              dvipdfmx math_note_template_v2.dvi
\documentclass[a5paper,10pt,onecolumn]{jsarticle}
\usepackage[dvipdfmx,top=22mm,bottom=20mm,inner=19mm,outer=18mm,headheight=13pt,headsep=6mm]{geometry}
\usepackage{amsmath,amssymb,amsthm,mathtools}
\usepackage{enumitem,fancyhdr,titlesec}
\usepackage[dvipdfmx,bookmarks=true,colorlinks=false]{hyperref}
\usepackage{pxjahyper}

% ----- 差し替える文書情報 -----
\newcommand{\NoteTitle}{凸双対性とLegendre--Fenchel変換}
\newcommand{\NoteSeries}{数学ノート集}

% ----- 文字・組版 -----
% 参考PDFの質素な明朝＋Computer Modernの印象を引き継ぎながら，
% A5判で窮屈にならないよう本文の行間を少し確保する．
\linespread{1.13}
\setlength{\parindent}{1zw}
\setlength{\parskip}{0pt}
\setlength{\abovedisplayskip}{6pt plus 2pt minus 2pt}
\setlength{\belowdisplayskip}{6pt plus 2pt minus 2pt}
\setlength{\abovedisplayshortskip}{3pt plus 2pt minus 1pt}
\setlength{\belowdisplayshortskip}{4pt plus 2pt minus 1pt}
\setlist[enumerate]{topsep=3pt,itemsep=1pt,parsep=0pt,leftmargin=2.4zw}
\titleformat{\section}{\normalfont\normalsize\bfseries}{\thesection}{0.8zw}{}
\titlespacing*{\section}{0pt}{2.0ex plus .4ex}{.9ex}
\titleformat{\subsection}{\normalfont\normalsize\bfseries}{\thesubsection}{0.7zw}{}
\titlespacing*{\subsection}{0pt}{1.5ex plus .3ex}{.6ex}
\raggedbottom

% ----- 柱とページ番号：ページ上部、罫線なし -----
\fancypagestyle{mathnote}{%
  \fancyhf{}%
  \fancyhead[L]{\footnotesize \NoteSeries\ （\NoteTitle）}%
  \fancyhead[R]{\footnotesize \thepage}%
  \renewcommand{\headrulewidth}{0pt}%
  \renewcommand{\footrulewidth}{0pt}%
}
\pagestyle{mathnote}

% ----- 節内通し番号の定義・定理群 -----
\newtheoremstyle{mathnotetheorem}
 {4.5pt}{4.5pt}{\normalfont}{}%
 {\bfseries}{}{0.65zw}%
 {\thmname{#1}\thmnumber{ #2}\thmnote{（#3）}}
\theoremstyle{mathnotetheorem}
\newtheorem{theorem}{定理}[section]
\newtheorem{definition}[theorem]{定義}
\newtheorem{lemma}[theorem]{補題}
\newtheorem{proposition}[theorem]{命題}
\newtheorem{corollary}[theorem]{系}
\newtheorem{example}[theorem]{例}
\newtheorem{remark}[theorem]{注意}
\renewcommand{\proofname}{証明}
\renewcommand{\qedsymbol}{$\square$}
\makeatletter
\renewenvironment{proof}[1][\proofname]{%
  \par\pushQED{\qed}\normalfont\topsep4\p@\@plus2\p@\relax
  \trivlist\item[\hskip\labelsep\bfseries #1]\ignorespaces
}{\popQED\endtrivlist\@endpefalse}
\makeatother

% ----- 使うノートで適宜拡張する記号 -----
\newcommand{\R}{\mathbb{R}}
\newcommand{\norm}[1]{\lVert #1\rVert}
\newcommand{\abs}[1]{\lvert #1\rvert}

\begin{document}
\section*{到達目標と前提}
凸共役と劣微分による最適性条件を証明し，
その条件から具体的な凸最適化問題の解を構成できるようになる．
双対問題は答えの正当性を検算する道具として用いる．
前提は実解析，連続線形汎関数，凸関数の初歩である．

\section{拡張実数値凸関数}

以下，$X$ を実Banach空間，$X^*$ を連続双対空間とし，
双対対合を $\langle x^*,x\rangle=x^*(x)$ と書く．

\begin{definition}
関数 $f:X\to(-\infty,+\infty]$ が
\emph{真}であるとは，恒等的に $+\infty$ ではなく，
$-\infty$ を値に取らないことをいう．
有効定義域を
$\operatorname{dom}f=\{x:f(x)<+\infty\}$ と書く．
$f$ が\emph{凸}であるとは
\[
 f((1-t)x+ty)\le(1-t)f(x)+tf(y)
\]
が $x,y\in X$，$0\le t\le1$ に対して成立することをいう．
$f$ が\emph{下半連続}であるとは
$\{x:f(x)\le a\}$ が全ての実数 $a$ について閉であることをいう．
\end{definition}

\begin{definition}
$f$ の\emph{エピグラフ}を
\[
 \operatorname{epi}f
 =\{(x,t)\in X\times\mathbb R:f(x)\le t\}
\]
と定義する．
\end{definition}

\begin{proposition}\label{prop:epi}
$f$ が凸であることと，
$\operatorname{epi}f$ が $X\times\mathbb R$ の凸集合であることは同値である．
また $f$ が下半連続であることと，
$\operatorname{epi}f$ が閉集合であることは同値である．
\end{proposition}
\begin{proof}
$(x,s),(y,t)\in\operatorname{epi}f$ なら
$f(x)\le s,f(y)\le t$ である．
凸性から
\[
 f((1-\lambda)x+\lambda y)
 \le(1-\lambda)f(x)+\lambda f(y)
 \le(1-\lambda)s+\lambda t
\]
となるのでエピグラフは凸である．
逆に $(x,f(x))$，$(y,f(y))$ に凸性を適用すれば，
$f$ の凸性を得る（有限値を持つ点について示せば十分）．

下半連続なら
$(x_n,t_n)\to(x,t)$ かつ $f(x_n)\le t_n$ に対し
$f(x)\le\liminf_nf(x_n)\le t$ であるから，
エピグラフは閉である．
逆にエピグラフが閉なら，任意の $a\in\mathbb R$ について
$\{x:f(x)\le a\}$ は
$\operatorname{epi}f\cap(X\times\{a\})$
の $X$ への同一視によって閉となり，下半連続である．
\end{proof}

\section{Legendre--Fenchel変換}

\begin{definition}
真な関数 $f:X\to(-\infty,+\infty]$ に対し，
\emph{凸共役関数} $f^*:X^*\to(-\infty,+\infty]$ を
\[
 f^*(x^*)=\sup_{x\in X}
 \{\langle x^*,x\rangle-f(x)\}
\]
によって定める．
\end{definition}

\begin{proposition}[凸共役関数の基本性質]\label{prop:conjugate-basic}
$f$ が真なら $f^*$ は凸かつ弱*下半連続である．
特にノルム下半連続でもある．
\end{proposition}
\begin{proof}
固定した $x\in\operatorname{dom}f$ に対し
$x^*\mapsto\langle x^*,x\rangle-f(x)$ は
弱*連続なアフィン関数である．
$f^*$ はこれらの関数の上限なので凸である．
任意の実数 $a$ について
\[
 \{x^*:f^*(x^*)\le a\}
 =\bigcap_{x\in\operatorname{dom}f}
 \{x^*:\langle x^*,x\rangle-f(x)\le a\}
\]
は弱*閉である．
従って $f^*$ は弱*下半連続である．
弱*位相はノルム位相より弱いので，
弱*閉集合はノルム閉でもある．
\end{proof}

\begin{theorem}[Fenchel--Youngの不等式]\label{thm:fy}
任意の $x\in X$，$x^*\in X^*$ について
\[
 f(x)+f^*(x^*)\ge\langle x^*,x\rangle
\]
が成り立つ．
\end{theorem}
\begin{proof}
$f^*(x^*)$ は
$\langle x^*,y\rangle-f(y)$（$y\in X$）
の上限なので，
特に $y=x$ として
$f^*(x^*)\ge\langle x^*,x\rangle-f(x)$ となる．
$f(x)=+\infty$ の場合は不等式は自明である．
\end{proof}

\section{劣微分と等号成立条件}

\begin{definition}
$f$ を真凸関数とする．$x\in\operatorname{dom}f$ に対し，
\[
 \partial f(x)=
 \{x^*\in X^*:f(y)\ge f(x)+\langle x^*,y-x\rangle
 \text{ for all }y\in X\}
\]
を $f$ の $x$ における\emph{劣微分}という．
\end{definition}

\begin{theorem}[Fenchel--Youngの等号]\label{thm:subgrad}
$f$ を真凸関数，$x\in\operatorname{dom}f$，
$x^*\in X^*$ とする．次は同値である．
\begin{enumerate}
\item $x^*\in\partial f(x)$．
\item $f(x)+f^*(x^*)=\langle x^*,x\rangle$．
\end{enumerate}
\end{theorem}
\begin{proof}
$x^*\in\partial f(x)$ なら
$f(y)\ge f(x)+\langle x^*,y-x\rangle$ なので
$\langle x^*,y\rangle-f(y)
 \le\langle x^*,x\rangle-f(x)$ が全ての $y$ で成立する．
$y=x$ で等号となるため，
上限を取って
$f^*(x^*)=\langle x^*,x\rangle-f(x)$ を得る．
逆にこの等号が成り立つなら，
任意の $y$ に対して
$\langle x^*,y\rangle-f(y)\le f^*(x^*)$
を用いれば
$f(y)\ge f(x)+\langle x^*,y-x\rangle$ となる．
従って $x^*\in\partial f(x)$ である．
\end{proof}

\begin{proposition}[内点での劣勾配の存在]\label{prop:subgrad-exists}
$f$ を真凸関数とし，$x\in\operatorname{int}(\operatorname{dom}f)$
で $f$ が $x$ の近傍で連続であるとする．
このとき $\partial f(x)\ne\varnothing$ である．
\end{proposition}
\begin{proof}
方向 $h\in X$ に対し
\[
 p(h)=\lim_{t\downarrow0}
 \frac{f(x+th)-f(x)}t
\]
と置く．$x$ は有効定義域の内点なので，
十分小さい正数 $t$ について分子は有限である．
凸関数の割線の傾きは $t$ について単調非減少であるから，
この極限は存在する．
$f$ が $x$ の近傍で連続だから，
ある $r>0$ について
$|f(x+u)-f(x)|\le1$ が $\norm u\le r$ で成り立つ．
これと凸性から $p(h)$ は有限であり，
$|p(h)|\le C\norm h$ となる定数 $C$ が存在する．

$p(th)=tp(h)$（$t\ge0$）は定義から従う．
また凸性から
\[
 f(x+t(h+k))
 \le\frac12f(x+2th)+\frac12f(x+2tk)
\]
となり，$t\downarrow0$ として
$p(h+k)\le p(h)+p(k)$ を得る．
従って $p$ は連続な劣線形汎関数である．
実Hahn--Banachの定理を零部分空間上の零汎関数に
適用すると，連続線形汎関数 $x^*$ で
$\langle x^*,h\rangle\le p(h)$
が全ての $h$ で成り立つものが存在する．
任意の $y\in\operatorname{dom}f$ について
凸性から
$p(y-x)\le f(y)-f(x)$ なので
$f(y)\ge f(x)+\langle x^*,y-x\rangle$ を得る．
$y\notin\operatorname{dom}f$ の場合は自明である．
従って $x^*\in\partial f(x)$ となる．
\end{proof}

\begin{corollary}[凸関数の最小値判定]\label{cor:optimality}
$f$ を真凸関数とし，$x\in\operatorname{dom}f$ とする．
$x$ が大域的最小点であることと
$0\in\partial f(x)$ は同値である．
\end{corollary}
\begin{proof}
$0\in\partial f(x)$ は
$f(y)\ge f(x)+0$ が全ての $y$ で成立することと同値であり，
これは $x$ が大域的最小点であることに他ならない．
\end{proof}

\section{Fenchel--Moreauの定理}

\begin{definition}
$f$ の\emph{双共役}を
\[
 f^{**}(x)=
 \sup_{x^*\in X^*}
 \{\langle x^*,x\rangle-f^*(x^*)\}
 \qquad(x\in X)
\]
と定める．
これは自然な埋め込み $X\hookrightarrow X^{**}$ を通じて
通常の二重共役を $X$ 上に制限したものである．
\end{definition}

\begin{lemma}[連続アフィン下界の存在]\label{lem:affine-minorant}
$f:X\to(-\infty,+\infty]$ を真・下半連続・凸関数とする．
このときある $x_0^*\in X^*$，$c\in\mathbb R$ が存在して
\[
 f(x)\ge\langle x_0^*,x\rangle+c
 \qquad(x\in X)
\]
となる．
\end{lemma}
\begin{proof}
$x_0\in\operatorname{dom}f$ を取り，
$a<f(x_0)$ とする．
$(x_0,a)$ は閉凸集合
$C=\operatorname{epi}f\subset X\times\mathbb R$ の外にある．
Hahn--Banachの強分離定理から，
$0$ でない連続線形汎関数
$\ell(x,t)=\langle y^*,x\rangle+\lambda t$ と
$\alpha\in\mathbb R$ が存在して
\[
 \ell(x,t)\le\alpha<\ell(x_0,a)
 \qquad((x,t)\in C)
\]
となる．
$C$ は $(x,t)\in C$ なら $(x,t+s)\in C$（$s\ge0$）
を満たすので $\lambda\le0$ である．
もし $\lambda=0$ なら，
$(x_0,f(x_0))\in C$ から
$\langle y^*,x_0\rangle\le\alpha$
となるが，右側の狭義不等式は
$\alpha<\langle y^*,x_0\rangle$ となり矛盾する．
従って $\lambda<0$ である．
$t=f(x)$ として
$\langle y^*,x\rangle+\lambda f(x)\le\alpha$
を変形すると
\[
 f(x)\ge\left\langle\frac{y^*}{-\lambda},x\right\rangle
             -\frac{\alpha}{-\lambda}.
\]
$f(x)=+\infty$ の点では自明だから，
この不等式が全域で成立する．
\end{proof}

\begin{theorem}[Fenchel--Moreauの定理]\label{thm:fenchel-moreau}
$f:X\to(-\infty,+\infty]$ を真・下半連続・凸関数とする．
このとき
\[
 f^{**}(x)=f(x)\qquad(x\in X)
\]
が成り立つ．
\end{theorem}
\begin{proof}
Fenchel--Youngの不等式から
$\langle x^*,x\rangle-f^*(x^*)\le f(x)$ なので
$f^{**}(x)\le f(x)$ が全ての $x$ で成り立つ．
逆向きを示すため，任意の $x_0\in X$ と
$a<f(x_0)$（$a\in\mathbb R$）を取る．
エピグラフ $C=\operatorname{epi}f$ は非空閉凸集合であり，
$(x_0,a)\notin C$ なので，強分離定理によって
\[
 \ell(x,t)=\langle y^*,x\rangle+\lambda t
 \le\alpha<\ell(x_0,a)
 \qquad((x,t)\in C)
 \tag{$*$}
\]
となる連続線形汎関数 $\ell$ と $\alpha$ が存在する．
$C$ の上方向の閉性から $\lambda\le0$ である．

$\lambda<0$ なら
$\ell(x,f(x))\le\alpha$ を変形し，
ある連続アフィン関数
$h(x)=\langle z^*,x\rangle+c$ で
$h\le f$ かつ $h(x_0)>a$ となるものが存在する．
実際
$z^*=y^*/(-\lambda)$，
$c=-\alpha/(-\lambda)$ と置けばよい．
このとき
$\langle z^*,x\rangle-f(x)\le-c$ が全ての $x$ で成立するので
$f^*(z^*)\le-c$ であり，
\[
 f^{**}(x_0)\ge\langle z^*,x_0\rangle-f^*(z^*)
 \ge h(x_0)>a
\]
となる．

$\lambda=0$ の場合は，
補題\ref{lem:affine-minorant}によって
既知のアフィン下界
$h_0(x)=\langle z_0^*,x\rangle+c_0\le f(x)$ を取る．
対応する汎関数
$\ell_0(x,t)=\langle z_0^*,x\rangle-t$
は $C$ 上で $\ell_0(x,t)\le-c_0$ を満たす．
不等式 $(*)$ の狭義の余裕は正なので，
十分小さい $\varepsilon>0$ に対して
\[
 \ell+\varepsilon\ell_0
\]
も $C$ 上の上界と $(x_0,a)$ との間で
狭義の分離を保つ．
実際，左辺の $C$ 上の上界は
$\alpha-\varepsilon c_0$，
右辺の値は
$\ell(x_0,a)+\varepsilon\ell_0(x_0,a)$ であり，
元の差 $\ell(x_0,a)-\alpha>0$ があるため
十分小さい $\varepsilon$ で大小関係は保たれる．
新しい汎関数の $t$ の係数は $-\varepsilon<0$ なので，
直前の場合に帰着し $f^{**}(x_0)>a$ を得る．

$a<f(x_0)$ は任意だったから
$f^{**}(x_0)\ge f(x_0)$ が成り立つ．
$f(x_0)=+\infty$ の場合にも，
任意の実数 $a$ について同じ議論を適用することで
$f^{**}(x_0)=+\infty$ を得る．
よって全ての $x_0$ で等号が成り立つ．
\end{proof}

\section{基本的な凸共役の計算}

\begin{example}[二次関数]\label{ex:quadratic}
$H$ を実Hilbert空間とし，
$f(x)=\frac12\norm x^2$ とする．
Rieszの表現定理によって $H^*$ を $H$ と同一視すると
\[
 f^*(y)=\frac12\norm y^2
\]
となる．
\end{example}
\begin{proof}
内積の平方完成によって
\[
 \langle y,x\rangle-\frac12\norm x^2
 =\frac12\norm y^2-\frac12\norm{x-y}^2
 \le\frac12\norm y^2.
\]
$x=y$ で等号となるので上限は $\norm y^2/2$ である．
\end{proof}

\begin{example}[ノルムの凸共役]\label{ex:norm-conjugate}
$f(x)=\norm x$ とすると
\[
 f^*(x^*)=
 \begin{cases}
 0&(\norm{x^*}\le1),\\
 +\infty&(\norm{x^*}>1).
 \end{cases}
\]
\end{example}
\begin{proof}
$\norm{x^*}\le1$ なら
$\langle x^*,x\rangle-\norm x\le0$ であり，
$x=0$ で零となる．
$\norm{x^*}>1$ なら，
ある $u\in X$，$\norm u=1$ について
$\langle x^*,u\rangle>1$ となる．
$x=tu$，$t>0$ とすると
\[
 \langle x^*,tu\rangle-\norm{tu}
 =t(\langle x^*,u\rangle-1)\longrightarrow+\infty
\]
なので共役値は無限大である．
\end{proof}

\begin{example}[指示関数と支持関数]\label{ex:indicator}
空でない集合 $C\subset X$ の指示関数を
\[
 \delta_C(x)=\begin{cases}0&(x\in C),\\+\infty&(x\notin C)\end{cases}
\]
とする．このとき
\[
 \delta_C^*(x^*)=\sup_{x\in C}\langle x^*,x\rangle
\]
となる．右辺を $C$ の支持関数という．
\end{example}
\begin{proof}
凸共役の定義において $x\notin C$ では
$\langle x^*,x\rangle-\delta_C(x)=-\infty$
となるため，上限に寄与するのは $x\in C$ のみである．
従って表示式を得る．
\end{proof}

\section{Fenchel双対性}

以下の節では $X,Y$ を有限次元実線形空間とする．
\begin{definition}
線形写像 $A:X\to Y$ と真凸関数
$f:X\to(-\infty,+\infty]$，
$g:Y\to(-\infty,+\infty]$ に対し，
主問題を
\[
 \inf_{x\in X}\{f(x)+g(Ax)\}
 \tag{P}
\]
とする．双対問題を
\[
 \sup_{y^*\in Y^*}
 \{-f^*(-A^*y^*)-g^*(y^*)\}
 \tag{D}
\]
と定める．
\end{definition}

\begin{theorem}[弱双対性]\label{thm:weak-duality}
任意の $x\in X$ と $y^*\in Y^*$ について
\[
 -f^*(-A^*y^*)-g^*(y^*)
 \le f(x)+g(Ax)
\]
が成り立つ．従って双対問題の上限は
主問題の下限を超えない．
\end{theorem}
\begin{proof}
Fenchel--Youngの不等式を $f$ に適用すると
\[
 f(x)+f^*(-A^*y^*)\ge-\langle y^*,Ax\rangle
\]
であり，$g$ について
\[
 g(Ax)+g^*(y^*)\ge\langle y^*,Ax\rangle
\]
である．両式を足せば
$f(x)+g(Ax)+f^*(-A^*y^*)+g^*(y^*)\ge0$
となり，移項して結論を得る．
\end{proof}

\begin{theorem}[強双対性の十分条件]\label{thm:strong-duality}
$X,Y$ を有限次元実線形空間とし，
$f:X\to(-\infty,+\infty]$，
$g:Y\to(-\infty,+\infty]$ を真・凸・下半連続関数，
$A:X\to Y$ を線形写像とする．
ある $x_0\in\operatorname{dom}f$ が存在して
$g$ が $Ax_0$ の近傍で有限かつ連続であると仮定する．
さらに主問題の下限 $v$ が有限であるとする．
このとき主問題と双対問題の最適値は一致し，
双対問題の上限は達成される．
\end{theorem}
\begin{proof}
摂動値関数
\[
 \varphi(z)=\inf_{x\in X}\{f(x)+g(Ax+z)\}
 \qquad(z\in Y)
\]
を定める．$f,g$ の凸性から $\varphi$ は凸である．
実際 $x_1,x_2$ を選び，両方の項に凸性を適用して
$z_1,z_2$ に関する不等式を得た後に下限を取ればよい．
$\varphi(0)=v\in\mathbb R$ である．
$g$ は $Ax_0$ の近傍で有限・連続なので，
ある近傍 $U$ において
\[
 \varphi(z)\le f(x_0)+g(Ax_0+z)<\infty
\]
となり，$\varphi$ は $0$ の近傍で上に有界である．
凸性と $\varphi(0)=v$ が有限であることから
$\varphi$ はこの近傍で $-\infty$ を取れない．
実際，もし $\varphi(z)=-\infty$ なら
$\varphi(0)\le(\varphi(z)+\varphi(-z))/2=-\infty$
となり矛盾する．
従って $\varphi$ は $0$ の近傍で有限となり，
有限次元空間上の有限値凸関数の局所連続性から
$0$ で連続である．
命題\ref{prop:subgrad-exists}を適用すると
$y^*\in\partial\varphi(0)$ が存在し，
\[
 \varphi(z)\ge\varphi(0)+\langle y^*,z\rangle
 =v+\langle y^*,z\rangle
\]
となる．
任意の $x\in X$ と $w\in Y$ について
$z=w-Ax$ と置くと
\[
 f(x)+g(w)\ge\varphi(w-Ax)
 \ge v+\langle y^*,w-Ax\rangle.
\]
従って
\[
 f(x)+\langle A^*y^*,x\rangle
 +g(w)-\langle y^*,w\rangle\ge v.
\]
$x$ と $w$ を独立に動かして下限を取ると
\[
 -f^*(-A^*y^*)-g^*(y^*)\ge v
\]
を得る．
弱双対性によって逆向きの不等式が成り立つので
等号であり，$y^*$ は双対問題の最大値を達成する．
\end{proof}

\begin{remark}
強双対性の成立は正則性条件に依存する．
本稿では有限次元で，$g$ がある実現可能点の像で
連続となる条件を用いた．
一般のBanach空間や相対内点による条件では，
位相的閉性などをより慎重に扱う必要がある．
\end{remark}

\section{最適性条件}

\begin{proposition}[双対等号の最適性条件]\label{prop:optimality}
主問題の実現可能点 $x$ と双対変数 $y^*$ について，
\[
 -A^*y^*\in\partial f(x),\qquad
 y^*\in\partial g(Ax)
\]
が成り立つなら，$x$ は主問題の最小点，
$y^*$ は双対問題の最大点であり，
両者の目的関数値は等しい．
\end{proposition}
\begin{proof}
定理\ref{thm:subgrad}により
\[
 f(x)+f^*(-A^*y^*)=-\langle y^*,Ax\rangle,
\]
\[
 g(Ax)+g^*(y^*)=\langle y^*,Ax\rangle
\]
が成り立つ．両式を足して整理すると
\[
 f(x)+g(Ax)
 =-f^*(-A^*y^*)-g^*(y^*)
\]
となる．弱双対性から，この等号が成立する
実現可能な $x,y^*$ はそれぞれ最適である．
\end{proof}

\begin{remark}
凸双対性は，元の関数を線形汎関数による評価から
再構成する理論である．
Fenchel--Moreauの定理はこの再構成を厳密に保証し，
強双対性は最適化問題において
原問題と双対問題の値が一致する条件を与える．
\end{remark}

\section{凸双対を用いて閾値処理を導く}
\begin{theorem}[絶対値正則化の厳密解]
$a\in\mathbb R$，$\lambda>0$ とし
\[
 F(x)=\frac12(x-a)^2+\lambda|x|
 \qquad(x\in\mathbb R)
\]
とおく．$F$ の最小点はただ一つで，
\[
 x_*=\operatorname{sgn}(a)(|a|-\lambda)_+
\]
に等しい（$a=0$ では右辺を0とする）．
この最小化問題に対する双対問題は
\[
 \max_{|y|\le\lambda}\left(ay-\frac12y^2\right)
\]
であり，最適値は原問題と一致する．
\end{theorem}
\begin{proof}
$|x|=\max_{|z|\le1}zx$ だから，
任意の $|y|\le\lambda$ について
\[
 F(x)\ge\frac12(x-a)^2+yx
       =\frac12(x-a+y)^2+ay-\frac12y^2.
\]
よって原問題の最小値は双対目的の最大値以上である．
双対目的は $-(y-a)^2/2+a^2/2$ なので，
唯一の最大点は
\[
 y_*=\begin{cases}
   -\lambda&(a<-\lambda),\\
   a&(|a|\le\lambda),\\
   \lambda&(a>\lambda).
 \end{cases}
\]
$x_*=a-y_*$ と置けば主張の表示式を得る．
$|a|\le\lambda$ では $x_*=0$ であり
$y_*x_*=\lambda|x_*|=0$．
$|a|>\lambda$ では $x_*$ と $y_*$ の符号が一致し，
$|y_*|=\lambda$ だから
$y_*x_*=\lambda|x_*|$．
従って上の二つの不等式はいずれも
$x=x_*$，$y=y_*$ で等号となり，
原問題と双対問題の最適値は一致する．
$F$ の二次項が狭義凸なので最小点は一意である．
\end{proof}
\begin{example}[二つの数値例]
$a=3$，$\lambda=1$ なら $y_*=1$，
$x_*=2$ で，最小値は
$F(2)=\frac12(2-3)^2+2=5/2$．
$a=2/5$，$\lambda=1$ なら
$y_*=2/5$，$x_*=0$ で最小値は $2/25$ である．
\end{example}
\section*{到達確認問題と解答}
\noindent\textbf{問題．}
$a=-4$，$\lambda=2$ について
原問題と双対問題の最適解および最適値を求めよ．

\medskip\noindent\textbf{解答．}
$y_*=-2$，$x_*=a-y_*=-2$．原問題の値は
$\frac12(-2+4)^2+2|-2|=2+4=6$，
双対の値は $(-4)(-2)-(-2)^2/2=8-2=6$ で一致する．

\end{document}
